\documentclass[10pt,a4paper]{amsart}
\usepackage[margin=19mm]{geometry}
\usepackage{amsmath,amssymb,mathtools}
\usepackage[scr=rsfs,cal=euler]{mathalfa}
\usepackage{xfrac}
\usepackage[unicode,hidelinks,bookmarks=false]{hyperref}
\newcommand{\ie}{i.e.\ }
\newcommand{\cf}{cf.\ }
\newcommand{\resp}{respectively\ }
\newcommand{\Subsection}{\subsection}
\providecommand{\nobreakdash}{\nobreak\hbox{-}}
\setlength{\emergencystretch}{2em}
\multlinegap0pt
\usepackage{bm}
\usepackage{bbm}
\usepackage{dsfont}
\usepackage{xspace}
\usepackage{cancel}
\usepackage{mathtools}
\usepackage{amsthm}
\makeatletter
\@ifundefined{theorem}{\newtheorem{theorem}{Theorem}}{}
\@ifundefined{example}{\newtheorem{example}[theorem]{Example}}{}
\@ifundefined{lemma}{\newtheorem{lemma}[theorem]{Lemma}}{}
\@ifundefined{proposition}{\newtheorem{proposition}[theorem]{Proposition}}{}
\makeatother
\title[Symmetrization of measures and the one-dimensional Poisson e]{Symmetrization of measures and the one-dimensional Poisson equation with Dirichlet boundary conditions}
\author{Christos Papadimitriou}
\date{Source: arXiv:2506.13311v1 (CC BY 4.0)}
\begin{document}
\maketitle

\noindent\emph{Source and licence.} Symmetrization of measures and the one-dimensional Poisson equation with Dirichlet boundary conditions. Version arXiv:2506.13311v1 (CC BY 4.0), distributed under the Creative Commons Attribution 4.0 International licence, as recorded in the arXiv licence metadata for this exact version and shown on the abstract page https://arxiv.org/abs/2506.13311v1. Full rights notice: licenses/arxiv-2506.13311-CC-BY-4.0.txt.\par\smallskip

\noindent\textit{Editorial scope.} Counted unit: the Lemma whose source label is \texttt{lemma karamata u f} in the article (source0.tex lines 582--605, source label \texttt{lemma karamata u f}), delivered together with the article's own complete original proof of it (source0.tex lines 606--650, one \texttt{proof} environment of 45 lines). No mathematics is written, repaired or re-derived: statement and proof are the article's own text line for line. Required context retained uncounted: paradeigma (lines 415--417); hardy polar (lines 429--436); karamata green (lines 559--569); 4.2 (lines 561--563); 4.3 (lines 565--567). Every definition, display and label of those lines is retained unchanged. Omitted from this package and not redistributed: 1--414, 418--428, 437--558, 564--564, 568--581, 651--995. Nothing inside a retained excerpt is omitted or altered: every retained body is delivered exactly as the article's lines are written, so no omission receipt of any kind accompanies this package. Citations: the retained excerpts carry no citation command at all, so no bibliography entry of the article is transcribed and no reference is redistributed. Reference adaptation: none was needed and none was made; every reference inside the delivered unit resolves inside it, so no unresolved reference or citation is produced. Printed slips kept literal: no printed word, symbol, hypothesis or number of the counted statement or of its proof was altered. Layout and class: the article's own class is \texttt{amsart} with packages including inputenc, amsmath, graphicx, svg, amsfonts, mathtools, amssymb, float, ragged2e, extarrows, dirtytalk, accents, upgreek, cancel; the wrapper is the lane's pinned \texttt{amsart} editorial wrapper at 10pt on A4, which restates the theorem-like environments the retained text needs and adds the article's own macro definitions that the retained text needs (none), with extra wrapper packages (bm, bbm, dsfont, xspace, cancel, mathtools). The delivered numbering is the unit's own. Assets: no retained span contains a figure environment, a graphics-inclusion command, a PDF-page inclusion, a table or a TikZ picture; no image file is copied into this package, the package declares no asset dependency and no image file is redistributed with it. Licence: version arXiv:2506.13311v1 of this article is distributed under the Creative Commons Attribution 4.0 International licence, as recorded in the arXiv licence metadata for this exact version and shown on the abstract page https://arxiv.org/abs/2506.13311v1. A full rights notice is delivered at licenses/arxiv-2506.13311-CC-BY-4.0.txt. Characters: every retained excerpt is pure ASCII, so the delivered engine, XeLaTeX with its default Latin Modern text font, reports no missing character.
\begin{example}\label{paradeigma}\textnormal{
Fix a point \(x_0\in[-\pi,\pi]\) and consider the Green's function \(G(x_0,\cdot):[-\pi,\pi]\to[0,+\infty)\) with one argument being constant. It is trivial to see that \(G(0,\cdot)\) is its own s.d.r and thus no polarization towards zero affects it. Assume that \(x_0\neq0\) and let \(b\in(-\pi,0)\cup(0,\pi)\). Let \(\bigl(G(x_0,\cdot)\bigr)_H\) be the polarization of \(G(x_0,\cdot)\) with respect to \(b\). If \(x_0<0\) (respectively \(x_0>0\)), then \(\bigl(G(x_0,\cdot)\bigr)_H\equiv G(x_0,\cdot)\) if and only if \(b\in(-\pi,x_0]\cup(0,\pi)\) (respectively \(b\in(-\pi,0)\cup[x_0,\pi)\)).}
\end{example}

\begin{proposition}\label{hardy polar}
Let \(0\leq f,g\in L^1[-\pi,\pi]\). Let \(b\in(-\pi,0)\cup(0,\pi)\) and \(H\) be the open half-line of \(\mathbb{R}\) that contains 0. Let \(f_H\) and \(g_H\) denote the polarizations of \(f\) and \(g\) with respect to \(b\). Then
\[\int_{-\pi}^\pi f(y)g(y)dy\leq\int_{-\pi}^\pi f_H(y)g_H(y)dy.\]
Equality holds if and only if the set \(A_H=B_H\cup C_H\) has zero Lebesgue measure, where
\[B_H=\{y\in H\cap[-\pi,\pi]:f(y)<f(2b-y)\text{ and }g(y)>g(2b-y)\}\]
and
\[C_H=\{y\in H\cap[-\pi,\pi]:f(y)>f(2b-y)\text{ and }g(y)<g(2b-y)\}.\]
\end{proposition}

\begin{lemma}\label{karamata green}
Let \(G\) be the Green's function for the Dirichlet problem on \([-\pi,\pi]\), let \(b\in(-\pi,0)\cup(0,\pi)\). If \(b>0\) (respectively \(b<0\)), let \(I\) be the interval \([b,\pi]\) (respectively \([-\pi,b]\)). For all \(x,y\in I\), the following two inequalities hold
\begin{equation}\label{4.2}
G(x,y)+G(x',y)\leq G(x,y')+G(x',y')
\end{equation}
and
\begin{equation}\label{4.3}
G(x,y)+G(x,y')\leq G(x',y)+G(x',y'),
\end{equation}
where \(x'\) and \(y'\) are the reflections of \(x\) and \(y\) with respect to \(b\), i.e. \(x'=2b-x\) and \(y'=2b-y\).\\
\end{lemma}

\begin{lemma}\label{lemma karamata u f}
Let \(0\leq f\in L^1[-\pi,\pi]\). Let \(b\in(-\pi,0)\cup(0,\pi),\) let \(I\) be as defined in Lemma \textnormal{\ref{karamata green}} and let \(f_H\) be the polarization of \(f\) with respect to \(b\). The following hold:
\begin{enumerate}
    \item[(a)] If \(x\in I\), let \(x'=2b-x\). Then
    \begin{equation}\label{uf syn uf}
    u_f(x)+u_f(x')\leq u_{f_H}(x)+u_{f_H}(x').
    \end{equation}
    \item[(b)] If \(x\in I\), let \(x'=2b-x\). Then
    \begin{equation}\label{uf uf polwsh polwsh}
    u_f(x)\leq u_{f_H}(x').
    \end{equation}
    \item[(c)] If \(b>0\) (respectively \(b<0\)), let \(x\in[-\pi,b]\) (respectively \(x\in [b,\pi]\)). Then 
    \begin{equation}\label{uf uf polwsh}
    u_f(x)\leq u_{f_H}(x).
    \end{equation}
    For \(b>0\) (respectively \(b<0\)), the following are equivalent:
    \begin{enumerate}
        \item[(i)] \((\ref{uf uf polwsh})\) holds as an equality for some \(x\in(-\pi,b]\) (resp. \(x\in [b,\pi)\)).
        \item[(ii)] \((\ref{uf uf polwsh})\) holds as an equality for all \(x\in(-\pi,b]\) (resp. \(x\in [b,\pi)\)).
        \item[(iii)] \(u_f\equiv u_{f_H}\) on \([-\pi,\pi]\).
        \item[(iv)] \(f=f_H\) almost everywhere on \([-\pi,\pi]\).
    \end{enumerate}
\end{enumerate}
\end{lemma}

\begin{proof}
Due to symmetry, it suffices to prove the lemma for \(b\in(-\pi,0)\).\\
(a) Let \(x\in[-\pi,b]\) and \(x'=2b-x\in[b,2b+\pi].\) From \((\ref{4.2})\) we have that \(G(x,\cdot)+G(x',\cdot)\) coincides with its polarization with respect to \(b\). Thus, by Proposition \ref{hardy polar},
    \begin{align*}
    u_f(x)+u_f(x')&=\int_{-\pi}^\pi\bigl(G(x,y)+G(x',y)\bigr)f(y)dy\\
    &\leq\int_{-\pi}^\pi\bigl(G(x,y)+G(x',y)\bigr)_H f_H(y)dy\\
    &=\int_{-\pi}^\pi\bigl(G(x,y)+G(x',y)\bigr)f_H(y)dy=
    u_{f_H}(x)+u_{f_H}(x').
    \end{align*}
 (b) Let \(x\in[-\pi,b]\) and consider \(x'=2b-x\in[b,2b+\pi].\) By the definition of polarization, \(f\equiv f_H\) a.e. on \((2b+\pi,\pi]\). 
It is easy to see that \(G(x,y)\leq G(x',y)\) for all \(y\in[2b+\pi,\pi]\), and thus\\
    \[\int_{2b+\pi}^\pi G(x,y)f(y)dy\leq\int_{2b+\pi}^\pi G(x',y)f(y)dy=\int_{2b+\pi}^\pi G(x',y)f_H(y)dy.\]
By the definition of polarization, we have \(f_H(y)=\min\{f(y),f(y')\}\) for all \(y\in[-\pi,b],\) where \(y'=2b-y\). Also, \(f_H(y')=\max\{f(y),f(y')\}\), since \(y'\in[b,2b-y]\), and thus \(f_H(y)\leq f_H(y')\) for all \(y\in[-\pi,b].\)\\
Note that \(\{f_H(y),f_H(y')\}=\{f(y),f(y')\}\), thus we can distinguish two possible cases:\\
either 
    \[f_H(y)=f(y)\text{ and } f_H(y')=f(y')\]
or
\[f_H(y)=f(y')\text{ and } f_H(y')=f(y).\]
Assume that the former case is true, then \(f(y)=f_H(y)\leq f_H(y')=f(y')\).
By Lemma \ref{karamata green}, we have \(G(x,y)-G(x',y)\leq G(x',y')-G(x,y'),\) and, thus
\[G(x,y)f(y)-G(x',y)f_H(y)\leq G(x',y')f_H(y')-G(x,y')f(y').\]
Similarly, if the latter case holds, then \(f(y')=f_H(y)\leq f_H(y')=f(y)\). Inequality (\ref{4.3}) implies that \(G(x,y')-G(x',y)\leq G(x',y')-G(x,y)\). Thus,
\[G(x,y')f(y')-G(x',y)f_H(y)\leq G(x',y')f_H(y')-G(x,y)f(y).\]
We conclude that in both cases, for \(y\in[-\pi,b]\) we have
\[G(x,y)f(y)+G(x,y')f(y')\leq G(x',y)f_H(y)+G(x',y')f_H(y'),\]
and, integrating over \([-\pi,b]\) with respect to \(y\), we get
\[\int_{-\pi}^b\bigl(G(x,y)f(y)+G(x,y')f(y')\bigr)dy\leq\!\!\int_{-\pi}^b\bigl(G(x',y)f_H(y)+G(x',y')f_H(y')\bigr)dy,\]
and with a change of variables this implies that
\[\int_{-\pi}^{2b+\pi} G(x,y)f(y)dy\leq\int_{-\pi}^{2b+\pi} G(x',y)f_H(y)dy.\]
We have already shown that
\[\int_{2b+\pi}^\pi G(x,y)f(y)dy\leq\int_{2b+\pi}^\pi G(x',y)f_H(y)dy,\]
and, adding these two inequalities, we end up with \((\ref{uf uf polwsh polwsh})\).\\
(c) Let \(x\in[b,\pi]\). By Example \ref{paradeigma}, \(G(x,\cdot)\equiv \bigl(G(x,\cdot)\bigr)_H\). Thus, by Proposition \ref{hardy polar}, we have
\begin{align*}
u_f(x)&=\int_{-\pi}^\pi G(x,y)f(y)dy\leq\int_{-\pi}^\pi \bigl(G(x,y)\bigr)_H f_H(y)dy\\
&=\int_{-\pi}^\pi G(x,y)f_H(y)dy=u_{f_H}(x).
\end{align*}
Considering the equivalence statement on Lemma \ref{lemma karamata u f}(c), it is trivial to see that \((iv)\implies(iii)\implies(ii)\implies(i)\), and to close the circle, we just have to show that \((i)\implies(iv)\).\\
Assume that \((\ref{uf uf polwsh})\) holds as an equality for some \(x\in[b,\pi).\)
By Proposition \ref{hardy polar}, we have that \(\lambda(B_H\cup C_H)=0\), where
\[B_H=\{y\in(b,\pi]:f(y)<f(2b-y)\text{ and }G(x,y)>G(x,2b-y)\}\]
and
\[C_H=\{y\in(b,\pi]:f(y)>f(2b-y)\text{ and }G(x,y)<G(x,2b-y)\}.\]
By Example \(\ref{paradeigma}\), \(G(x,y)> G(x,2b-y)\) for all \(y\in(b,2b+\pi)\). Note that we have considered  \(G\) to  be identically zero outside of \([-\pi,\pi]^2\). Thus, \(C_H=\varnothing\) and \(B_H=\{y\in(b,2b+\pi):f(y)<f(2b-y)\}\). So, \(f=f_H\) a.e. on \((b,2b+\pi)\) and, by the definition of polarization, \(f=f_H\) a.e. on \([-\pi,\pi]\).
\end{proof}

\end{document}
